\section{Introduction}
\label{sec:intro}

\paragraph{What the paper proposes.}
A quantum neural network is a small quantum circuit whose rotation angles are treated as
trainable weights; the network's output is the probability of measuring one designated qubit in
the state $\ket{1}$.  Rosenhahn and Hirche \cite{ptqnn} propose a particular building block,
the \emph{prototype neuron}, in which each input qubit is first rotated by a learned angle, then
used to control a rotation of the neuron's own qubit, and finally rotated back
(\cref{fig:circuits}).  The effect is that the neuron ``fires'' most strongly when the input
pattern matches a stored prototype pattern, and the fact that the inputs are rotated back means
the same inputs can be shared by many neurons and by many training examples \emph{in
superposition}.  The paper demonstrates this on four tasks: learning a target histogram over all
32 five-bit strings at once (bit-count), classifying the Iris flowers, controlling a Flappy Bird
game, and denoising an electrocardiogram (ECG).  Training uses no gradient descent; it uses a
finite-difference estimate of the slope inside a simulated-annealing loop (\cref{sec:kwsa}).

\paragraph{What we did.}
This document records three things.  First, a \textbf{reproduction}: a clean-room implementation
(no code from the authors), cross-checked against an independent simulator written blind by a
second agent, that regenerates every figure and table of the paper that can be regenerated from
its text (\cref{sec:reproduction}).  Second, a \textbf{noise study}: a description of the seven
quantum computers for which vendors publish calibration data, as mathematical noise channels, and
an exact simulation of each experiment under each of them (\cref{sec:noise}), plus a
\textbf{compilation and cost study} of what it would take to run the circuits on each machine
(\cref{sec:hardware}).  Third, a \textbf{shot-noise programme}: having found that the number of
repeated measurements, not gate error, is what limits a hardware demonstration, we implement and
test three remedies (\cref{sec:shots}).

\paragraph{How to read the numbers.}
Every quantitative statement below comes from a \emph{run directory} --- a folder under
\code{experiments/runs/} holding the script's output together with the random seed, the commit
hash of the code and the simulator backend that produced it.  Run directories are named by
time stamp (for example \run{bitcount\_noise/20260820T145335Z}) and are quoted in the text; the
figures are produced from them by one script (\code{experiments/report\_figures.py}).  Where the
paper leaves a choice open (its optimiser step sizes, its train/test split, its game physics) we
had to choose, and each such choice is listed in \cref{app:assumptions}.  Every calibration
number of a quantum computer carries, in the repository, the URL it was read from, the date, and
a verbatim quotation (\cref{app:provenance}).

\paragraph{Compute discipline.}
Two machines were used.  Run directories from 25 September 2026 record the machine (host name,
platform, cores used) in their \code{meta.json}; for earlier ones the machine is attributed by
date.  The August 2026 runs --- the
reproduction, the noise study and the hardware plan, i.e.\ everything whose time stamp begins
\code{202608} --- ran on one workstation (64 cores, one RTX~3090 GPU).  Three runs of
18 September 2026 (\run{bitcount/20260918T111515Z}, \run{ecg/20260918T132924},
\run{iris/20260918T085641Z-drop}) ran on a 12-thread laptop with no GPU; the runs of
25 September (\code{20260925}) ran on the workstation again, from clean commits, and reproduce
to the last recorded digit every laptop number they overlap (the laptop run directories are not
on the workstation; the comparison is against the values recorded from them).  Wall times are quoted with the number of cores
actually used.  No single run was allowed to exceed about two
hours; two studies that would have (noise-aware training of the 11-qubit bit-count network,
$\approx 13$ hours for its grid; the full ECG/Iris noise grid, $\approx 8$ hours) were
deliberately \emph{not} run and are reported as such, with their measured per-cell cost, rather
than approximated.  The ECG reruns on the paper's own recording use the $1\,800$-sample
Fig.~13 window, with the full recording as a check (\cref{sec:ecg}).

\begin{table}[t]
\centering\small
\caption{Scoreboard: what the paper reports and what we obtain.  ``Reproduced'' means the
qualitative claim and the quantitative value (within the precision the paper's figure allows)
both come out; ``discrepancy'' means they do not, and the text explains what we tried.}
\label{tab:scoreboard}
\begin{tabularx}{\textwidth}{@{}>{\raggedright\arraybackslash}p{0.21\textwidth} >{\raggedright\arraybackslash}p{0.21\textwidth} X >{\raggedright\arraybackslash}p{0.16\textwidth}@{}}
\toprule
Paper item & Paper value & Ours & Status \\
\midrule
Fig.~3, activation surface & $(0,1)\!\to\!\max$, $(1,0)\!\to\!0$ & $1.000$ / $0.000$ / $0.500$ / $0.500$ at the corners & reproduced \\
Fig.~5, target histogram & bar heights & $(\text{number of ones}+1)/112$, all 32 bars within $0.0002$ & reproduced \\
Fig.~6, superposition training & $\KL\approx 0$ in $\approx 8$ iterations & median $<0.005$ at sweep 9; $5/6$ seeds $<10^{-3}$ at 20 & reproduced \\
Fig.~6, iterative baseline & $\KL\approx 0$ by $\approx 120$ & per-input $\KL<10^{-3}$ for $4/6$ seeds within $10^5$ elementary steps (a fifth reaches $1.4\times10^{-3}$); the superposition $\KL$ is a different observable & reproduced in the per-input observable; iteration unit assumed \\
Fig.~7, Iris learning curves & classic overfits, prototype does not & classic test loss rises after 1--4 sweeps; prototype's keeps falling; over 20 seeds the classic test loss ends $>5\,\%$ above its minimum on 8 (merge) / 7 (drop), the prototype's on 2 / 0 & reproduced \\
Iris accuracy & $\approx 98\,\%$ & $95$ / $91$ / $98\,\%$ on three splits (ceilings $97$ / $99$ / $99\,\%$); over 20 seeds $94.3\pm2.8\,\%$ (classic $93.9\pm2.7\,\%$), $96.4\,\%$ on the filtered set when ambiguous codes are dropped & $2$--$4$ points short; no architecture advantage \\
Fig.~9, Iris weight picture & prototype weights span the data & t-SNE: both weight sets cluster away from the data.  Scored over 20 seeds: the prototype layer covers more codes than the classic one, but already does so untrained; training makes it more selective, yet less so than the classic layer; angles do not decode to data patterns & discrepancy (criterion fixed in advance) \\
Flappy Bird & ``plays perfectly'' & $110/110$ pipes (the cap) on $20/20$ held-out games & reproduced (own clone) \\
Table~I, ECG PSNR (1/2/3 hidden) & $32.9$ / $40.00$ / $40.04$~dB & $37.4$ / $38.0$ / $38.1$~dB (synthetic signal); $36.8$ / $37.5$ / $36.9$~dB on the paper's own recording ($37.8$ / $37.8$ / $36.2$~dB on all $7\,700$ samples) & within $\approx 2$~dB on synthetic; on the original recording 1-neuron reproduced, 2/3-neuron $2$--$4$~dB short \\
\bottomrule
\end{tabularx}
\end{table}
